% Part: turing-machines
% Chapter: undecidability
% Section: trakhtenbrot

\documentclass[../../../include/open-logic-section]{subfiles}

\begin{document}

\olfileid{tur}{und}{tra}
\olsection{ट्राख्टेनब्रोटचे प्रमेय}

\begin{explain}
  \olref[rep]{sec} मध्ये ट्यूरिंग
  यंत्र~$M$ आणि आदान
  चिन्हमाला~$w$ यांच्यासाठी
  $!T(M,w)$ आणि~$!E(M,w)$
  ही !!{sentence}s वाक्ये
  परिभाषित केली. त्यानंतर
  \olref[ver]{lem:valid-if-halt}
  आणि \olref[ver]{lem:halt-if-valid}
  मध्ये आदान~$w$ वर
  सुरू केलेले~$M$ अखेरीस
  थांबते तेव्हाच
  $!T(M,w) \lif !E(M,w)$
  वैध असते, हे दाखवले.
  थांबण्याची समस्या अनिर्णेय
  असल्यामुळे प्रथम-क्रम
  तर्कशास्त्रातील !!{sentence}s
  वाक्यांची वैधता आणि
  पूर्ततायोग्यताही अनिर्णेय आहेत
  (\cref{tur:und:uns:thm:decision-prob,tur:und:uns:cor:undecidable-sat}).

  पण वाक्यांची वैधता आणि
  पूर्ततायोग्यता कोणत्याही
  !!{structure}s रचनांसाठी
  परिभाषित होतात, मग त्या
  सांत असोत वा अनंत.
  एखादे !!a{sentence} वाक्य
  सांत !!{structure} रचनेत
  पूर्ततायोग्य आहे का
  (किंवा सर्व सांत
  !!{structure}s रचनांमध्ये
  वैध आहे का), हे ठरवणे
  सोपे असेल असे वाटू शकते.
  निर्णय समस्या सोडवता
  येत नाही याचा पुरावा
  आपण यासाठीही रूपांतरित
  करू शकतो; तसे नसते
  हे त्यातून दिसेल.

  प्रथम \olref[ver]{lem:halt-if-valid}
  च्या पुराव्याकडे परत पाहा.
  तेथे आपण~$!T(M,w)$ चे
  प्रतिरूप~$\Struct M$
  तयार केले; आदान~$w$
  वर सुरू केलेले यंत्र~$M$
  नेमके काय करते ते
  त्यात वर्णिले आहे.
  त्या प्रतिरूपाचे परिभाषाक्षेत्र
  $\Nat$, म्हणजे अनंत,
  होते. पण $M$ हे
  आदान~$w$ वर खरोखर
  थांबत असेल, तर त्याच
  पद्धतीने सांत प्रतिरूप
  ~$\Struct M'$ बांधता येते.
  समजा आदान~$w$ वर
  सुरू झालेले $M$ हे
  $k$ पायऱ्यांनंतर थांबते.
  $\Domain{M'}$ हे
  $\{0, \dots, n\}$ घ्या;
  येथे $n$ ही संख्या
  $k$ पेक्षा एकाने मोठी
  संख्या आणि~$w$ च्या
  लांबीपेक्षा एकाने मोठी
  संख्या यांपैकी मोठी
  आहे. तसेच
  \[
    \Assign{\prime}{M'}(x) = 
    \begin{cases}
      x + 1 &\text{$x < n$ असेल तर}\\
      n &\text{अन्यथा,}
    \end{cases}
  \]
  घ्या, आणि
  $\tuple{x,y} \in \Assign{<}{M'}$
  तेव्हाच, जेव्हा $x < y$
  किंवा $x = y = n$.
  इतर बाबतीत $\Struct{M'}$
  ची व्याख्या~$\Struct{M}$
  प्रमाणेच आहे. $\Struct{M'}$
  च्या व्याख्येनुसार,
  \olref[ver]{lem:halt-if-valid}
  च्या पुराव्यातल्याप्रमाणे
  $\Sat{M'}{!T(M,w)}$.
  तसेच $M$ हे आदान~$w$
  वर थांबते असे गृहीत
  धरल्यामुळे
  $\Sat{M'}{!E(M,w)}$.
  म्हणून~$\Struct{M'}$
  हे $!T(M,w) \land !E(M,w)$
  चे सांत प्रतिरूप आहे
  (येथे $\lif$ ऐवजी
  $\land$ वापरले आहे).
%READERNOTE{OLTUR-017 — गोठवलेल्या इंग्रजी सांत प्रतिरूपात वरची मर्यादा थांबण्याच्या पायऱ्या k आणि आदानाची लांबी यांतील मोठी संख्या घेतली आहे. यंत्र घर एकपासून सुरू होऊन k उजव्या हालचालींनंतर k+1 या घरात पोहोचू शकते. तसेच सर्वोच्च घटकावर n<n ठेवलेल्या सांत क्रमामुळे n हे शेवटचे आदान-चिन्ह असलेले घर नसावे; अन्यथा रिकाम्या-शेपटीची अट त्या घरालाच रिकामे म्हणेल. म्हणून मराठीत दोन्ही संख्यांपेक्षा एकाने मोठी मर्यादा घेतली आहे.}
  
  पुराव्याचा अर्धा भाग
  झाला आहे: $M$ हे
  आदान~$w$ वर थांबले,
  तर $!T(M,w) \land !E(M,w)$
  ला सांत प्रतिरूप
  असते, हे दाखवले.
  दुर्दैवाने याचा उलट
  निष्कर्ष खरा नाही:
  काही ट्यूरिंग यंत्रे
  आदान~$w$ वर थांबत
  नाहीत, तरी
  $!T(M,w)
  \land !E(M,w)$ ला
  सांत प्रतिरूप असू शकते.
  उदाहरणार्थ, एकच
  अवस्था $q_0$ आणि
  $\delta(q_0,\TMblank) = \tuple{q_0,\TMblank,\TMstay}$
  ही एकच सूचना
  असलेले यंत्र~$M$
  घ्या. रिकाम्या आदानावर
  $w = \emptyseq$
  सुरू केल्यावर ते
  कधीच थांबत नाही:
  ते अनंत फेरीत
  राहते, पण फीत
  बदलत नाही किंवा
  शीर्ष हलवत नाही.
  सर्व स्थितिवर्णने
  सारखीच असतात
  (तीच अवस्था,
  तेच शीर्षस्थान,
  तीच फीत).
  तरी $!T(M,\emptyseq) \land !E(M,\emptyseq)$
  सत्य करणारी सांत
  !!{structure} रचना
  ~$\Struct{M''}$
  परिभाषित करता येते
  (सराव). पण अशा
  !!{structure}s रचना
  वगळण्यासाठी $!T(M,w)$
  योग्य रीतीने बदलता
  येते.

  \begin{prob}
    $q_0$ ही एकच
    अवस्था आणि
    $\delta(q_0,\TMblank) = \tuple{q_0,\TMblank,\TMstay}$
    ही एकच सूचना
    असलेले ट्यूरिंग
    यंत्र~$M$ घ्या.
    $\Domain{M''} = \{0, 1, 2\}$,
    $\Assign{\prime}{M''}(0) =
    \Assign{\prime}{M''}(1) = 1$
    आणि $\Assign{\prime}{M''}(2) = 2$,
    तसेच $\Assign{<}{M''} =
    \{\tuple{0,1}, \tuple{1,1}, \tuple{2,2}\}$
    असे घ्या.
    $\Assign{\Obj Q_{q_0}}{M''}$,
    $\Assign{\Obj
    S_{\TMblank}}{M''}$
    आणि $\Assign{\Obj S_{\TMendtape}}{M''}$
    अशी परिभाषित करा
    की $!T(M,\emptyseq)$
    आणि $!E(M,\emptyseq)$
    दोन्ही सत्य ठरतील;
    तसे का होते ते
    स्पष्ट करा.
    सूचना: $\delta(q_0, \TMendtape)$
    अपरिभाषित आहे हे
    लक्षात घ्या.
    पुढील दोन्ही विधाने
    ~$\Struct{M''}$
    मध्ये सत्य आहेत
    याची खात्री करा:
    \begin{align*}
    & \Obj Q_{q_0}(\num{1}, \num{n}) \land \Obj S_{\TMendtape}(\num{0}, \num{n})
      \land 
      \lforall[x][(\num{0} < x \lif \Obj S_\TMblank(x, \num{n}))] \text{\quad सर्व $n \in \Nat$ साठी}\\
    & \lexists[y][(\Obj Q_{q_0}(\num{0}, y) \land \Obj S_{\TMendtape}(\num{0}, y))]
    \end{align*}
  \end{prob}
%READERNOTE{OLTUR-018 — गोठवलेल्या इंग्रजी सरावातील एकमेव सूचना पुढील अवस्थेला अव्याख्यायित q म्हणते, जरी एकच अवस्था q0 दिली आहे; तसेच M'' च्या उत्तरवर्ती फलनात एका ठिकाणी M' लिहिले आहे. मराठीत दोन्ही ठिकाणी दिलेलीच एक-अवस्थेची यंत्रणा आणि M'' रचना वापरली आहे.}
\end{explain}

ट्यूरिंग यंत्र~$M$ हे
आदान $w = \sigma_{i_1}\dots\sigma_{i_k}$
वर कसे चालते, याचे
वर्णन करणारी पुढील
!!{sentence}s वाक्ये पाहा:
\begin{enumerate}
\item संख्या आणि~$<$
यांचे वर्णन करणारी
स्वयंसिद्धके (\olref[rep]{sec}
मधील~$!T(M,w)$ च्या
व्याख्येप्रमाणेच).
\item आरंभीच्या स्थितिवर्णनाचे
वर्णन करणारी स्वयंसिद्धके:
तीही~$!T(M,w)$
च्या व्याख्येप्रमाणेच.
\item एका स्थितिवर्णनातून
पुढच्या स्थितिवर्णनात
होणाऱ्या संक्रमणाचे
वर्णन करणारी स्वयंसिद्धके:

पुढील बाबींसाठी $!A(x, y)$
आधीप्रमाणे असू द्या,
आणि
\[
  !B(y) \ident \lforall[x][(x < y \lif \eq/[x][y])].
\]
\begin{enumerate}
\item \ollabel{rep-right}
$\delta(q_i, \sigma) =
  \tuple{q_j, \sigma', \TMright}$
रूपाच्या प्रत्येक सूचनेसाठी
पुढील !!{sentence} वाक्य:
\begin{align*}
& \lforall[x][\lforall[y][(
   (\Obj Q_{q_i}(x, y) \land \Obj S_{\sigma}(x, y)) \lif {}]] \\
&\qquad   (\Obj Q_{q_j}(x', y') \land \Obj S_{\sigma'}(x, y') \land
!A(x, y) \land !B(y')))
\end{align*}
\item \ollabel{rep-left}
$\delta(q_i, \sigma) =
  \tuple{q_j, \sigma', \TMleft}$
रूपाच्या प्रत्येक सूचनेसाठी
पुढील !!{sentence} वाक्य:
\begin{align*}
& \lforall[x][\lforall[y][
    ((\Obj Q_{q_i}(x', y) \land \Obj S_{\sigma}(x', y)) \lif {}]]\\
& \qquad   (\Obj Q_{q_j}(x, y') \land \Obj S_{\sigma'}(x', y') \land
!A(x', y) \land !B(y'))) \land {}\\
& \lforall[y][((\Obj Q_{q_i}(\num{0}, y) \land \Obj S_{\sigma}(\num{0},
    y)) \lif {}]\\
& \qquad (\Obj Q_{q_j}(\num{0}, y') \land \Obj S_{\sigma'}(\num{0},
  y') \land !A(\num{0}, y) \land !B(y')))
\end{align*}
\item \ollabel{rep-stay}
$\delta(q_i, \sigma) =
  \tuple{q_j, \sigma', \TMstay}$
रूपाच्या प्रत्येक सूचनेसाठी
पुढील !!{sentence} वाक्य:
\begin{align*}
& \lforall[x][\lforall[y][(
   (\Obj Q_{q_i}(x, y) \land \Obj S_{\sigma}(x, y)) \lif {}]] \\
&\qquad (\Obj Q_{q_j}(x, y') \land \Obj S_{\sigma'}(x, y') \land
!A(x, y) \land !B(y')))
\end{align*}
\end{enumerate}
येथे दिसते त्याप्रमाणे~$M$
च्या संक्रमणांचे वर्णन
करणारी !!{sentence}s
वाक्ये~$!T(M,w)$ मधील
संबंधित !!{sentence}
वाक्यांसारखीच आहेत;
फक्त शेवटी $!B(y')$
जोडले आहे. $!B(y')$
मुळे “पुढच्या” स्थितिवर्णनाचा
क्रमांक $y'$ हा आधीच्या
सर्व क्रमांकांहून वेगळा
असतो: $\Obj 0$,
$\Obj 0'$, \dots.
%READERNOTE{OLTUR-019 — गोठवलेल्या इंग्रजी सांत-प्रतिरूप बांधणीत डावीकडे जाण्याच्या सकारात्मक-स्थान शाखेत स्थिरता-अट लिहिलेल्या घराऐवजी नव्या शीर्षस्थानाला वगळते; त्याच शाखेत नव्या वेळेसाठी B अटही गाळली आहे, जरी लगतचा मजकूर प्रत्येक संक्रमणात ती जोडल्याचे सांगतो. येथे आधीच्या OLTUR-009 नुसार लिहिलेले घर वगळले आहे आणि या शाखेतही नव्या वेळेची अट घातली आहे. शून्यावरील शाखा बदललेली नाही.}
% \item सुसंगततेच्या अटी व्यक्त करणारी वाक्ये:
% \begin{enumerate}
%   \item\ollabel{rep-con-symb}%
%   एका वेळी एका घरात एकापेक्षा अधिक
%   चिन्हे असू शकत नाहीत असे सांगणारे !!^a{sentence}:
%   \[\lforall[x][\lforall[y][(\Obj S_{\sigma}(x, y) \lif \lnot \Obj
%   S_{\sigma'}(x, y))]],\]
%   $T$ च्या फितीवरील प्रत्येक भिन्न चिन्ह-जोडी $\sigma \neq \sigma'$ साठी.
%   \item\ollabel{rep-con-state}%
%   एका वेळी $M$ एकापेक्षा अधिक
%   अवस्थांत असू शकत नाही असे सांगणारे !!^a{sentence}:
%   \[\lforall[x][\lforall[y][(\Obj Q_{q}(x, y) \lif
%   \lforall[z][\lnot \Obj
%   Q_{q'}(z, y))]],\]
%   $T$ यंत्राच्या प्रत्येक भिन्न अवस्था-जोडी $q \neq q'$ साठी.
%   \item\ollabel{rep-con-tape}%
%   एका वेळी $M$ फितीवरील एकापेक्षा अधिक
%   घरांवर असू शकत नाही असे सांगणारे !!^a{sentence}:
%   \[\lforall[x][\lforall[y][(\Obj Q_{q}(x, y) \lif
%   \lforall[z][\eq/[z][y]\lif \lnot \Obj
%   Q_{q'}(z, y))]],\]
%   $T$ यंत्राच्या प्रत्येक अवस्था-जोडी $q$, $q'$ साठी.
% \end{enumerate}
\end{enumerate}
$!T'(M, w)$ हे ट्यूरिंग
यंत्र~$M$ आणि आदान~$w$
यांच्यासाठी वरील सर्व
!!{sentence}s वाक्यांचे
संयुक्त विधान असू द्या.

\begin{lem}\ollabel{lem:halts-sat}
  आदान~$w$ वर सुरू
  केलेले $M$ थांबत असेल,
  तर $!T'(M,w) \land !E(M,w)$
  ला सांत प्रतिरूप असते.
\end{lem}

\begin{proof}
  \olref[ver]{lem:halt-if-valid}
  च्या पुराव्यातील
  $\Struct{M'}$ घ्या;
  मात्र पुढील बदल करा:
  \begin{align*}
    \Domain{M'} & = \{0, \dots, n\},\\
    \Assign{\prime}{M'}(x) & = 
    \begin{cases}
      x + 1 &\text{$x < n$ असेल तर}\\
      n &\text{अन्यथा,}
    \end{cases} \\
    \tuple{x,y} \in \Assign{<}{M'} &\text{$x < y$ किंवा $x = y = n$ असेल तेव्हाच,}
  \end{align*}
  येथे $n = \max(k+1,\len{w}+1)$
  आणि~$k$ ही अशी
  लघुतम संख्या आहे
  की आदान~$w$ वर
  सुरू झालेले $M$
  हे~$k$ पायऱ्यांनंतर
  थांबले आहे.
  $\Sat{M'}{!T'(M,w) \land !E(M,w)}$
  याची पडताळणी सरावासाठी
  सोडतो.
%READERNOTE{OLTUR-020 — गोठवलेल्या इंग्रजी लेम्माच्या पुराव्यात आणि लगेचच्या सरावात थांबण्याचे वाक्य E असे लिहिले आहे; सरावात T' ऐवजी जुने T आले आहे. येथे दोन्ही ठिकाणी लेम्मातील संपूर्ण T' आणि !E संयुक्त विधानच तपासायचे आहे.}
\end{proof}

\begin{prob}
  $\Sat{M'}{!T'(M,w) \land !E(M,w)}$
  हे सिद्ध करून
  \olref[tur][und][tra]{lem:halts-sat}
  चा पुरावा पूर्ण करा.
\end{prob}

\begin{lem}\ollabel{lem:sat-halts}
  $!T'(M,w) \land !E(M,w)$
  ला सांत प्रतिरूप
  असेल, तर आदान~$w$
  वर सुरू झालेले
  $M$ थांबते.
\end{lem}

\begin{proof}
  प्रतिपरिवर्तन सिद्ध करू.
  समजा आदान~$w$
  वर सुरू झालेले
  $M$ थांबत नाही.
  $!T'(M,w) \land !E(M,w)$
  ला प्रतिरूपच नसेल,
  तर पुरावा पूर्ण.
  म्हणून~$\Struct{M}$
  हे~$!T'(M,w) \land !E(M, w)$
  चे प्रतिरूप आहे
  असे समजा. ते
  सांत असू शकत
  नाही हे दाखवू.
%READERNOTE{OLTUR-021 — गोठवलेल्या इंग्रजी प्रतिपरिवर्तनात गृहीत प्रतिरूप अचानक जुन्या T आणि !E संयुक्त विधानाचे म्हटले आहे, जरी लेम्मा आणि पुढील निष्पन्नता T' वर अवलंबून आहेत. येथे लेम्मातीलच T' आणि !E गृहीत घेतले आहेत.}
  
  \olref[ver]{lem:config}
  प्रमाणे, आदान~$w$
  वर सुरू केलेले~$M$
  हे~$n$ पायऱ्यांनंतर
  थांबलेले नसेल आणि
  किमान एक पायरी
  झालेली असेल, तर
  $!T'(M,w)
  \Entails !C(M, w, n) \land !B(\num{n})$
  हे सिद्ध करता येते.
  आदान~$w$ वर
  सुरू झालेले~$M$
  थांबत नसल्यामुळे,
  प्रत्येक सकारात्मक
  $n \in \Nat$ साठी
  $!T'(M,w) \Entails !C(M, w, n) \land
  !B(\num{n})$.
  \olref[rep]{prop:mlessk}
  नुसार सर्व~$k < n$
  साठी $!T'(M,w) \Entails \num{k} < \num{n}$
  हेही लक्षात घ्या.
  तसेच $!B(\num{n})
  \Entails \num{k} < \num{n} \lif
  \eq/[\num{k}][\num{n}]$.
  म्हणून सर्व~$k < n$
  साठी $\Sat{M}{\eq/[\num{k}][\num{n}]}$;
  म्हणजे अनंत संख्यांक
  पदे~$\num{k}$ यांना
  $\Struct{M}$ मध्ये
  परस्परभिन्न मूल्ये
  असलीच पाहिजेत.
  त्यासाठी $\Domain{M}$
  अनंत असावे लागते;
  म्हणून~$\Struct{M}$
  हे~$!T'(M,w) \land !E(M, w)$
  चे सांत प्रतिरूप
  असू शकत नाही.
%READERNOTE{OLTUR-022 — गोठवलेल्या इंग्रजी मजकुरात आरंभीच्या वेळेसाठीही B(0) निष्पन्न असल्याचा दावा आहे. पण B(y') ही अट संक्रमणांच्या निष्कर्षातच आहे; आरंभीच्या वेळेसाठी ती दिलेली नाही. येथे B(n) फक्त किमान एक पायरी झाल्यावर वापरले आहे. परस्परभिन्न संख्यांक मिळवण्यासाठी हेच पुरेसे आहे.}
\end{proof}

\begin{prob}
  आदान~$w$ वर
  सुरू केलेले~$M$
  हे~$n$ पायऱ्यांनंतर
  थांबलेले नसेल आणि
  किमान एक पायरी
  झालेली असेल,
  तर $!T'(M,w) \Entails !B(\num{n})$
  हे सिद्ध करून
  \olref[tur][und][tra]{lem:sat-halts}
  चा पुरावा पूर्ण करा.
\end{prob}

\begin{thm}[ट्राख्टेनब्रोटचे प्रमेय]
  \ollabel{thm:trakhtenbrodt}
  प्रथम-क्रम तर्कशास्त्रातील
  कोणत्याही !!{sentence}
  वाक्याला सांत प्रतिरूप
  आहे का (म्हणजे ते
  सांत प्रतिरूपात पूर्ततायोग्य आहे का),
  हे ठरवणे अनिर्णेय आहे.
\end{thm}

\begin{proof}
  सांत प्रतिरूपात पूर्ततायोग्यता ठरवणारे
  ट्यूरिंग यंत्र~$F$
  आहे असे समजा.
  मग कोणतेही ट्यूरिंग
  यंत्र~$M$ आणि
  आदान~$w$ दिले असता
  $!T'(M,w) \land !E(M,w)$
  हे वाक्य संगणित करून
  त्याला सांत प्रतिरूप
  आहे का हे~$F$
  ने ठरवता येईल.
  \cref{tur:und:tra:lem:halts-sat,tur:und:tra:lem:sat-halts}
  नुसार त्याला सांत
  प्रतिरूप असते तेव्हाच
  आदान~$w$ वर
  सुरू झालेले $M$
  थांबते. म्हणून~$F$
  वापरून थांबण्याची
  समस्या सोडवता येईल;
  पण ती सोडवता येत
  नाही हे आपल्याला
  माहीत आहे.
\end{proof}

\begin{cor}\ollabel{cor:fproof-incomp}%
  \emph{सांत} वैधतेसाठी
  निर्दोष आणि संपूर्ण
  अशी कोणतीही !!{derivation}
  पद्धती असू शकत नाही.
  म्हणजे प्रत्येक सांत
  !!{structure} रचना~$\Struct{M}$
  साठी $\Sat{M}{!B}$
  असेल तेव्हाच
  $\Proves !B$ असेल,
  अशी !!a{derivation}
  पद्धती असू शकत नाही.
\end{cor}

\begin{proof}
  सराव.
\end{proof}

\begin{prob}
  \olref[tur][und][tra]{cor:fproof-incomp}
  सिद्ध करा. प्रत्येक
  सांत !!{structure}
  रचनेत $!B$ सत्य
  असेल तेव्हाच
  $\lnot !B$ सांत प्रतिरूपात
  पूर्ततायोग्य नसेल,
  हे लक्षात घ्या.
  \olref[tur][und][uns]{thm:valid-ce}
  मधील अर्थाने सांत
  प्रतिरूपात पूर्ततायोग्यता अर्धनिर्णेय
  का आहे, हे स्पष्ट
  करा. सांत वैधतेसाठी
  !!a{derivation} पद्धती
  असेल, तर सांत
  प्रतिरूपात पूर्ततायोग्यता निर्णेय
  होईल, असा युक्तिवाद
  करण्यासाठी हे वापरा.
\end{prob}


\end{document}
